\subsection{All anabelian even bases, including permuted simple factors}

\begin{statement}{D.44}{Proved}{The entire relative pronilpotent family over anabelian even bases and its exact solvable-top residual}
Call a finite group anabelian here when all its composition factors are
nonabelian simple. This definition includes the trivial group. With
the original marked sign map, put
\[
 \mathscr A=\{N\lhd B:N\subset B_{\rm even},\
                         B_{\rm even}/N\text{ is finite anabelian}\}.
                                                               \tag{D.44.1}
\]
For \(N\in\mathscr A\), write \(J_N=B/N\) and use the actual
relative completions and source kernels
\[
 \begin{split}
 R_{N,p}&=\ker(\widehat N\longrightarrow N^{(p)}),\\
 \mathcal P_N&=\widehat B/\bigcap_pR_{N,p},\\
 R_{\mathscr A}&=\bigcap_{N\in\mathscr A,\ p}R_{N,p},\\
 \mathcal P_{\mathscr A}&=\widehat B/R_{\mathscr A},\\
 J_{\mathscr A}&=\varprojlim_{N\in\mathscr A}J_N .
 \end{split}                                                   \tag{D.44.2}
\]
The family is directed under intersections, with all maps induced
by the same \(B\). Its actual simultaneous completion and base
kernel are
\[
 \mathcal P_{\mathscr A}=\varprojlim_N\mathcal P_N,\qquad
 \mathcal N_{\mathscr A}:=\ker(\mathcal P_{\mathscr A}\to J_{\mathscr A})
                  =\prod_p\varprojlim_N N^{(p)}.                \tag{D.44.3}
\]
Every displayed transition map on the pro-\(p\) kernels is
surjective. Thus this kernel is pronilpotent and projects onto
the kernel at every finite base and every finite relative level.
The maximal prosolvable quotient of \(\mathcal P_{\mathscr A}\)
is exactly \(\mathcal R_3\) of D.42.

This includes every marked semisimple even base, each with its
entire relative pronilpotent kernel, and every finite joint base
formed from them. In an individual base having a transitive
permutation action on its simple factors, that orbit has size
one or two. The two-factor source (D.43.11) is included with its
actual swapping involution. The larger family also retains
successive anabelian even extensions; no factor-permutation
action is discarded.

Keep \(K_D,\mathcal D,\mathcal W,K_{T,C},\mathcal C_{54}\)
and the actual maps from D.42--D.43. Define
\[
 \begin{split}
 \mathcal D_{\mathscr A}
       &=\widehat B/(K_D R_{\mathscr A}),\\
 \mathcal C_{\mathscr A}
       &=\mathcal C_{54}\times_{C_6}\mathcal D_{\mathscr A},\\
 \mathcal Z_{44}
       &=\widehat B/(R_W\cap R_{\mathscr A})
         =\mathcal W\times_{\mathcal C_{\mathscr A}}
                              \mathcal P_{\mathscr A}.
 \end{split}                                                   \tag{D.44.4}
\]
Here \(\mathcal C_{\mathscr A}\) is the exact maximal common
marked quotient, not only a quotient of the common image.
There is an actual continuous peripheral automorphism
\[
            (\alpha_{41},1_{\mathcal P_{\mathscr A}})
                   \quad\hbox{of }\mathcal Z_{44}.              \tag{D.44.5}
\]
Its multiplier is one. It retains the D.39 action and every
D.40--D.43 coordinate with their actual maps and common gauge.
One simultaneous pro-\(3\) fixed-element lift supplies its
loop conjugators. If \(\delta:\widehat B\to\mathcal D\)
is the original source map, the exact projection kernels are
\[
 \begin{split}
 \ker(\mathcal Z_{44}\to\mathcal P_{\mathscr A})
       &=K_{T,C}\times\delta(R_{\mathscr A}),\\
 \ker(\mathcal Z_{44}\to\mathcal W)
       &=\ker(\mathcal P_{\mathscr A}\to\mathcal C_{54})
           \cap
         \ker(\mathcal P_{\mathscr A}\to\mathcal D_{\mathscr A}).
 \end{split}                                                   \tag{D.44.6}
\]

There is a smaller solvable even quotient at the next unresolved
source. Let \(a=(1\,2)(3\,4)\), \(b=(1\,3\,5)\) in \(A_5\),
and use the wreath-product convention
\[
 (v,q)(w,h)=\bigl((v_iw_{q^{-1}i})_{i=1}^4,qh\bigr).
\]
The marked homomorphism
\[
 \begin{split}
 \rho:\widehat B&\longrightarrow E=A_5^4\rtimes S_4,\\
 s&\longmapsto((1,1,a,1),(1\,2)),\\
 r&\longmapsto((b,1,1,1),(2\,3\,4))
 \end{split}                                                   \tag{D.44.7}
\]
is onto. Its whole image has order \(311040000\); its even
image has order \(155520000\), and its simple base
\(M=A_5^4\) has order \(12960000\).
The actual orders of \(t\) and \(x\) are \(20\) and \(10\).
The top map is precisely the canonical
\(B\to\mathrm{SL}_2(\mathbb Z/4)/\{\pm I\}\simeq S_4\),
with the indicated images of \(s,r\). The remaining source kernel
has the exact image
\[
              \rho(R_W\cap R_{\mathscr A})=A_5^4.              \tag{D.44.8}
\]
The even quotient above this kernel is \(A_4\), with derived
orders \(4,1\). This exhibits an actual nonnilpotent solvable
even quotient beyond the anabelian family. D.45 tests its two
minimal order-\(24\) relative pro-\(2\) bases and proves
nonextension of the recorded action. Neither (D.44.5) nor
(D.44.8) decides full realization of its entire prescribed CM pair.

These comparisons address Part D(e), section 51.16,
folio 734 verso, scan 174, and D-after-LM2(1),
section 51.15, folios 725--727, scans 156--160.
\end{statement}
\begin{proof}
\emph{Directed marked bases and the full relative kernels.}
Normal subgroups and quotients of a finite anabelian group are
anabelian, by a composition series. A finite subdirect product
of such groups is anabelian: for two factors, its projection
kernel is a normal subgroup of the other factor, and induction
handles any finite number. Applied to the actual even image
of \(B\), this proves that \(\mathscr A\) is directed under
intersections. The actual element \(s\) splits every \(J_N\)
over its even subgroup. Since that subgroup has only nonabelian
composition factors, it is perfect and centreless.
Every \(N\) with nontrivial even base is torsion-free: the order-two factor survives
the sign, and killing \(r\) would make the nontrivial even
image generated by \(r,srs\) trivial. For the trivial even
base \(N=B_{\rm even}=C_3*C_3\), the pro-\(p\) completion
is used literally, without the free-rank formula. All other
\(N\) are free of rank \(1+[B:N]/6\), by their marked
Bass--Serre graphs.

If \(N_1\subset N_2\), the finite quotient \(N_2/N_1\)
is a normal subgroup of the even group \(B_{\rm even}/N_1\)
and is perfect. In every finite \(p\)-group image of \(N_2\),
the image of \(N_1\) is the whole image: the quotient would
otherwise be a nontrivial \(p\)-group quotient of a perfect
group. Hence
\[
                    N_1^{(p)}\twoheadrightarrow N_2^{(p)}.    \tag{D.44.9}
\]
The inclusion of the actual source kernels gives the
corresponding relative quotient maps. The finite relative
quotients over all finite bases form an inverse system with
the original marked source \(B\). Compactness proves that
its limit is exactly the quotient by the intersection in
(D.44.2). Exactness of inverse limits of these surjective
compact systems, followed by reordering the prime
coordinates, proves (D.44.3). In particular its kernel
projects onto every finite-base kernel, a property needed
below.

In a finite relative image let \(V\) denote the image of
the nilpotent kernel. Its even quotient modulo \(V\) is
anabelian and perfect. After the perfect residual of the
whole image is killed, the remaining even subgroup is an
image of \(V\). It is nilpotent and generated by the two
actual order-three elements \(r,srs\). It is consequently
a \(3\)-group. Conversely the exact relative pro-\(3\)
source over the sign is a quotient of every relative
completion: \(N^{(3)}\to B_{\rm even}^{(3)}\) is onto by
the same perfect-quotient argument. This proves the
maximal prosolvable assertion, with its actual marked map.

For a semisimple even base, conjugation by its even
subgroup preserves each simple factor. Thus the
permutation action of the whole base on these factors
factors through its sign \(C_2\). A transitive orbit
therefore has size one or two. All these bases, their
finite subdirect joint images, and their relative
completions belong to the directed system just proved.
In particular the actual \(A_7^2\) swapping source in
D.43 is retained. This argument uses its permutation
action; it does not replace it by two separately
normalized simple coordinates.

\emph{The common quotient, with both comparison maps.}
We give the precise extension of the finite common-image
argument of D.43. Let \(Q_*\) be any finite common quotient
of \(\mathcal W\) and \(\mathcal P_{\mathscr A}\).
Let \(P=Q_*^\infty\), the image of the properfect
factor \(\mathcal R'\times\prod A_n\), and let \(Z\)
be the image of the isolated coordinate
\(\mathcal T'\to\mathcal W,\ z\mapsto(z,1,1)\).
Then \(Z\) is normal and centralizes \(P\).
Factor through one finite relative base using the
directed system, and let \(V\) be the image of its
nilpotent kernel. The even subgroup of \(Q_*/V\)
is an anabelian quotient, hence is centreless and
perfect. The image of \(P\) surjects onto it.
The even image of \(Z\), which centralizes that
subgroup, is trivial. Therefore \(Z\subset V\).

For \(p\ne3\), its characteristic Sylow \(Z_p\)
has trivial image in \(Q_*/P\), whose even subgroup
is a \(3\)-group. Thus \(Z_p\subset Z(P)\).
For any even \(q\in Q_*\), choose \(a\in P\)
having the same image modulo \(V\); conjugation
on \(Z_p\) is then that of \(a^{-1}q\in V\).
This action is a \(p\)-group, whereas the even
source is generated by the two order-three
elements \(r,srs\). The action is trivial.
The actual invariant map
\[
 B'\longrightarrow\mathcal T'\longrightarrow Z_p
\]
therefore has normal kernel in \(B_{\rm even}\).
It yields a central extension of
\(B_{\rm even}/B'=C_3\), generated by those two
order-three elements. Such an extension is
abelian of exponent dividing three. Its
\(p\)-subgroup \(Z_p\) vanishes. Thus \(Z\)
is a \(3\)-group. This uses a map from \(B'\);
it introduces no projection from all of
\(B_{\rm even}\) onto \(\mathcal T'\).

The same actual \(\mathcal T\)-quotient as in
D.43 now factors through \(\mathcal C_{54}\).
Every common quotient kills the isolated
\(K_{T,C}\). Put \(K_C=\ker(\widehat B\to
\mathcal C_{54})\). Since
\(R_{\mathscr A}\subset K_C\), the two
inclusions just proved and the modular
identity for normal subgroups give
\[
 R_W R_{\mathscr A}
   =(K_C\cap K_D)R_{\mathscr A}
   =K_C\cap(K_D R_{\mathscr A}).                               \tag{D.44.10}
\]
All products are closed, being images of
products of compact groups. The maximal
solvable quotient of \(\mathcal D_{\mathscr A}\)
is \(C_6\), since its commutator is an image
of the properfect commutator of \(\mathcal D\).
Both factors retain the literal \(C_6\).
Their maximal common quotient is therefore
\(C_6\). Closed Goursat applied to
(D.44.10) proves both directions of
(D.44.4). Taking its projection kernels
gives (D.44.6).

\emph{One simultaneous normalized lift.}
The recorded \(\alpha_{41}\) is the identity
on \(\mathcal C_{54}\) and \(\mathcal D\).
Its action on their common quotient is
therefore the identity, so the product
map (D.44.5) is an automorphism.
Its recorded conjugators have common images
\[
       A_C=1,\qquad \beta_C=(b_C,1),\qquad
                    b_C=-\tfrac12(e_7+e_8).                   \tag{D.44.11}
\]
In every finite common image, the preceding
proof places the isolated \(\mathcal T'\)
image in the image of the nilpotent base
kernel and makes it a \(3\)-group.
The surjectivity in (D.44.3) and compactness
therefore place \(\beta_C\) in the image of
\[
                    \mathcal N_{\mathscr A,3}
                             =\varprojlim_N N^{(3)}.
\]
The map onto that image is equivariant
for \(s\), with pro-\(3\) kernel.
The actual element \(\beta_C\) is
\(s\)-fixed. Coprime fixed-element lifting
provides one \(s\)-fixed
\(\widetilde\beta\in\mathcal N_{\mathscr A,3}\)
above it. For completeness, take finite
quotients normal in
\(\mathcal N_{\mathscr A,3}\rtimes\langle s\rangle\).
Each finite system has a fixed lift by
coprime centre induction; finite
intersections are handled by their common
refinement. Compactness gives one lift
at every base, prime and homology level
simultaneously.

The pairs \((A_W,1)\) and
\((\beta_W,\widetilde\beta)\) are consequently
actual elements of (D.44.4). They give
the required conjugations of \(r,s\).
The automorphism fixes \(t\), and its
parameter is \((f_{39},1)\). Its actual
source preimage lies in \(\widehat F'\),
because the \(\mathcal W\) coordinate
retains the full abelianization of
\(\widehat F\). Both hexagons follow
from the same order-two, order-three
and fixed-cusp equations as in D.25.
The gauge and all projected actions
are unchanged. This establishes a
quotient automorphism; D.45 supplies
an explicit obstruction to extending
this recorded action to the full source.

\emph{The marked wreath source and its exact residual image.}
The fixed-block choices in (D.44.7)
give \(s^2=r^3=1\). The top generators
generate \(S_4\). The program
\texttt{d\_cm\_anabelian\_even.py}
enumerates the \(24\) top elements
and all \(48\) Schreier generators
of its kernel. Their first-coordinate
images generate all \(60\) elements
of \(A_5\). Top conjugation gives
surjectivity to each coordinate.
A subdirect product of four simple
groups is a product of diagonals
over a partition of the coordinates,
by the simple Goursat argument.
Normality makes that partition
\(S_4\)-invariant. Its primitive
action leaves only the single-block
and singleton partitions.
The certificate records a literal
\(22\)-letter \(B\)-word whose
first coordinate is \(1\) and
whose other coordinates are nontrivial.
This excludes a full diagonal and
proves that the kernel is \(A_5^4\).
The group orders follow.
The program also checks all \(576\)
products for the canonical mod-four
top isomorphism and the exact orders
of the displayed cusp elements.

The subgroup \(M=A_5^4\) is the unique
minimal normal subgroup of \(E\).
Indeed its factors form one top orbit;
any normal subgroup meeting \(M\)
nontrivially contains \(M\).
Moreover \(C_E(M)=1\): a centralizer
cannot permute factors, and each
simple factor has trivial centre.
A normal subgroup disjoint from
\(M\) centralizes \(M\), so is
trivial. Thus every nontrivial
normal subgroup of \(E\) contains \(M\).

There is no quotient map
\(\mathcal W\twoheadrightarrow E\).
Such a map would factor through
a finite simultaneous image of
its solvable coordinate and
finitely many perfect and symmetric
coordinates. It cannot factor
through any one coordinate:
\(E\) is nonsolvable and is not
perfect; and it is no quotient
of any \(S_n,\ n\ge5\).
The latter's nonabelian quotient
is \(S_n\) itself, whose derived
subgroup is simple, whereas
\(E'\) is not simple.
Each coordinate kernel would
have a nontrivial normal image
containing \(M\). Their iterated
commutator is trivial in the
subdirect source but contains
the perfect \(M\) in \(E\),
a contradiction.

There is likewise no quotient
\(\mathcal P_{\mathscr A}\twoheadrightarrow E\).
The image of the pronilpotent base
kernel would be a nilpotent
normal subgroup. Every nontrivial
normal subgroup contains the
perfect \(M\), so that image
is trivial. The map would then
factor through an anabelian
even base. But the even image
of \(E\) has the nontrivial
solvable quotient \(A_4\).
This is impossible by its
composition factors.
Finally, an epimorphism from
their actual joint image
\(\mathcal Z_{44}\) would give
nontrivial images of both
coordinate kernels, each
containing \(M\). Those kernels
commute in the fibre product;
the perfect nontrivial \(M\)
cannot commute with itself.
Thus \(E\) is not a quotient
of \(\mathcal Z_{44}\).

The image \(\rho(R_W\cap
R_{\mathscr A})\) is consequently
a nontrivial normal subgroup
and contains \(M\). It is
contained in \(M\), because
\(\mathcal W\) already retains
the exact canonical \(S_4\)
top through \(Q\) modulo four.
This proves (D.44.8), with
both inclusions and the literal
marked source map.
\end{proof}
